\documentclass[11pt]{article}

\usepackage[utf8]{inputenc}
\usepackage[T1]{fontenc}
\usepackage{lmodern}
\usepackage{geometry}
\usepackage{amsmath,amssymb,amsfonts,amsthm,mathtools}
\usepackage{bm}
\usepackage{enumitem}
\usepackage{microtype}
\usepackage{hyperref}
\usepackage{titlesec}
\usepackage{booktabs}
\usepackage{array}
\usepackage{setspace}
\usepackage{mathrsfs}
\usepackage{physics}

\geometry{margin=1in}
\hypersetup{colorlinks=true, linkcolor=black, urlcolor=blue, citecolor=black}
\setlist[enumerate]{topsep=0.4em,itemsep=0.35em}
\setlist[itemize]{topsep=0.3em,itemsep=0.25em}
\titleformat{\section}{\large\bfseries}{\thesection.}{0.5em}{}
\titleformat{\subsection}{\normalsize\bfseries}{\thesubsection.}{0.5em}{}
\setstretch{1.06}

\newcommand{\Aether}{\AE{}ther}
\newcommand{\Aflow}{\AE{}ther-flow}
\newcommand{\TheoryName}{\Aflow{} Interpretation of Relativity}
\newcommand{\TheoryProgram}{\Aether{} / \Aflow{} framework}
\newcommand{\CompletedMetric}{g^{\sharp}}

\newcommand{\PullbackBodyResponseCore}{\mathfrak S^{\mathrm{sub,proto,pppresec,pb,bodyresp}}}
\newcommand{\SubstrateLaw}{\mathfrak L^{\mathrm{sub,disp,resp}}}

\newtheorem{definition}{Definition}
\newtheorem{proposition}{Proposition}
\newtheorem{remark}{Remark}
\newtheorem{corollary}{Corollary}
\newtheorem{theorem}{Theorem}

\title{\textbf{The \TheoryName{}}\\[0.4em]
\large Primitive-Reservoir Observer-Reduction\\
\large Section-Centered Hidden-Response\\
\large Pullback Body-Response Substrate-Law\\
\large Next Benchmark Criterion}
\author{}
\date{}

\begin{document}

\maketitle

\begin{abstract}
The current active-primary theorem now derives the current body-response core from a benchmark-faithful section-centered displacement-response substrate law \(\SubstrateLaw\). That core is therefore no longer the live unexplained stopping point on the recorded section-centered hidden-response route. Another same-output deeper relay that merely preserves that same substrate law and therefore merely reproduces the same downstream one-metric observer package no longer counts as the honest default next benchmark-facing move.

This note records the routing consequence of that gain. The next honest benchmark-facing manuscript must now derive or force the current displacement-response substrate law itself, or some nontrivial constituent or relation inside it, from still more primitive explicit substrate structure in a way that changes benchmark-facing status. The claim boundary remains fixed: the one operative metric is still exactly \(\CompletedMetric\), the ordinary matter route is unchanged, there is no second operative metric, the low-energy matter law is not enlarged, and no stronger promotion language is licensed. The positive result is therefore routing discipline below the new body-response-substrate-law frontier.
\end{abstract}

\section{Introduction}

The active derivational program of \TheoryName{} has again moved the live burden upstream. The previous active-primary theorem proved that the observer-relevant pullback body-response core \(\PullbackBodyResponseCore\) is the unique minimal benchmark-facing datum on the recorded section-centered hidden-response route. The new theorem goes beyond that internal sharpening: it shows that the same current body-response core is itself forced by a benchmark-faithful section-centered displacement-response substrate law \(\SubstrateLaw\).

That changes what now counts as honest primary progress. Once the body-response core is no longer primitive, another manuscript that merely reproduces the same core by replaying the same substrate-law data does not remove a remaining benchmark-facing burden. The next honest move must therefore target the substrate law itself, or a genuinely smaller theorem inside it.

\paragraph{Framework and claim boundary.}
\TheoryProgram{} denotes the broader \Aether{} / \Aflow{} framework built on the claims that the \Aether{} is the underlying four-dimensional substrate of reality, the \Aflow{} is its intrinsic ordered motion, observed three-dimensional space is the local experiential slice of that deeper substrate, S-time is the experienced order of change arising from matter, light, and the \Aflow{}, and observed expansion is the three-dimensional appearance of deeper four-dimensional ordered motion. Within that framework, \TheoryName{} denotes the exact-closure theory in which the effective gravitational dynamics are adopted to be exactly Einsteinian with universal matter coupling. In the active sequence, \emph{adoption} means use of that established relativistic dynamics without claiming substrate derivation, while \emph{derivation} is reserved for a first-principles recovery from explicit substrate variables.

The present note stays strictly inside that boundary. It records only which kind of next manuscript now counts as an honest benchmark-facing gain below the new substrate-law frontier.

\section{Fixed Inputs}

\begin{definition}[Substrate-law frontier inputs]
The criterion recorded here uses the following fixed inputs.
\begin{enumerate}
    \item The benchmark exact-closure package, which fixes the public one-metric GR target of the repository.
    \item The benchmark gatekeeping layer, including the recorded safeguard that blocks same-output deeper-origin relays from counting as new front-line progress once the live benchmark-relevant datum is unchanged.
    \item The earlier theorem proving that the current body-response core \(\PullbackBodyResponseCore\) is the unique minimal benchmark-facing datum on the recorded route up to body-response-faithful equivalence.
    \item The later theorem deriving that same current body-response core from a benchmark-faithful section-centered displacement-response substrate law \(\SubstrateLaw\).
\end{enumerate}
\end{definition}

\begin{remark}
The present note does not replace those manuscripts. It records the routing consequence of reading them together under the active safeguard against recursive depth drift.
\end{remark}

\section{Why the Live Burden Now Sits Below the Substrate-Law Frontier}

\begin{proposition}[The live burden is renewed below substrate-law restatement]
At the current stage of the primitive-reservoir observer-reduction line, the body-response-substrate-law theorem renews the live benchmark-facing burden below any mere restatement of the current section-centered displacement-response substrate law \(\SubstrateLaw\).
\end{proposition}

\begin{proof}
The new theorem changes benchmark-facing status because it removes the current body-response core \(\PullbackBodyResponseCore\) itself as the live unexplained stopping point. Once that core is forced by \(\SubstrateLaw\), another manuscript that merely preserves the same substrate law no longer removes any remaining benchmark-facing freedom. The live burden therefore no longer sits on restating the same law. It sits on deriving or sharpening that already-active law from still more primitive explicit substrate structure.
\end{proof}

\begin{definition}[Benchmark-neutral same-law relay below the substrate-law frontier]
A \emph{benchmark-neutral same-law relay below the substrate-law frontier} is a manuscript that preserves the same benchmark one-metric output, the same ordinary matter route, the same current section-centered displacement-response substrate law \(\SubstrateLaw\) up to substrate-law-faithful equivalence, the same current observer-relevant pullback body-response core \(\PullbackBodyResponseCore\), and the same downstream benchmark-facing consequences while merely relocating the unexplained substrate layer deeper, renaming the same law, or re-expanding back to the larger pullback-body-operator, pullback-operator, pullback-quadratic, hidden-response, residual-normal-form, or pre-proto-section presentations without a new derivation, new forcing theorem, or comparably sharp new gain.
\end{definition}

\section{The Current Post-Substrate-Law Criterion}

\begin{definition}[Admissible next benchmark-facing manuscript below the substrate-law frontier]
Below the current section-centered displacement-response substrate-law frontier, a manuscript counts as an admissible next benchmark-facing move only if it does at least one of the following:
\begin{enumerate}
    \item derives or justifies the current displacement-response substrate law \(\SubstrateLaw\) itself from still more primitive explicit substrate structure in a genuinely benchmark-facing way;
    \item derives or justifies some nontrivial constituent of \(\SubstrateLaw\), such as the substrate body, the displacement projection / minimal-lift splitting, or the coercive substrate-response operator, from still more primitive explicit substrate structure in a way that sharpens the present route;
    \item proves a sharper necessity, uniqueness, or compatibility theorem for some nontrivial constituent or relation inside \(\SubstrateLaw\);
    \item does one of the previous three while preserving the same benchmark one-metric package, the same ordinary matter route, and the same claim boundary.
\end{enumerate}
\end{definition}

\begin{theorem}[Next honest benchmark-facing task below the substrate-law frontier]
The next honest benchmark-facing manuscript below the current section-centered displacement-response substrate-law frontier must derive or justify the current law \(\SubstrateLaw\), or some nontrivial constituent or relation inside it, from still more primitive explicit substrate structure in a benchmark-relevant way, or else prove a genuinely sharper theorem inside that same law. A manuscript that only repeats the same current law through another deeper relay, re-expands the discussion back to the larger pullback-body-operator, pullback-operator, pullback-quadratic, hidden-response, residual-normal-form, or pre-proto-section presentations without such a new gain, or invokes a side branch without doing the same benchmark-facing work does not satisfy the criterion.
\end{theorem}

\begin{proof}
By Proposition 1, the live burden already lies below any mere restatement of the current substrate law \(\SubstrateLaw\). Therefore a subsequent manuscript must improve on that already-active law rather than merely accompany it. A deeper-origin manuscript can do so only if it changes benchmark-facing status by more than depth alone. If it simply reproduces the same law and its same downstream outputs, the routing safeguard classifies it as benchmark neutral. The remaining honest alternatives are therefore exactly those stated in the definition: derive or justify the law itself, derive or justify some nontrivial constituent or relation inside it, or prove a sharper internal theorem there.
\end{proof}

\section{What the Next Move Should Not Be}

\begin{proposition}[Screened next moves below the substrate-law frontier]
Below the current section-centered displacement-response substrate-law frontier, none of the following is the default next benchmark-facing move:
\begin{enumerate}
    \item another same-output deeper-origin relay that merely preserves the current displacement-response substrate law \(\SubstrateLaw\) up to substrate-law-faithful equivalence and therefore merely reproduces the same current body-response core \(\PullbackBodyResponseCore\) and the same downstream benchmark-facing consequences;
    \item another re-expansion back to the larger pullback-body-operator, pullback-operator, pullback-quadratic, hidden-response, residual-normal-form, or pre-proto-section presentations as if those already-spent larger presentations were again independently live burdens;
    \item a side branch that does not derive or justify the same substrate law, derive or justify a nontrivial constituent or relation inside it, or close a benchmark derivation gate in a genuinely new way;
    \item a manuscript that introduces a second operative metric, enlarges the ordinary matter law, or uses stronger promotion language.
\end{enumerate}
\end{proposition}

\begin{proof}
Item (1) fails the preceding criterion because it leaves the renewed live burden unchanged. Item (2) likewise fails because it re-expands to larger already-spent presentations rather than deriving or sharpening the already-active substrate law. Item (3) does not directly address the live burden at current scope unless it performs the same benchmark-facing work named above. Item (4) violates the fixed claim boundary of the benchmark package and therefore cannot count as a disciplined next step on the active line. The benchmark package remains the exact one-metric GR package already fixed by adoption, so any admissible next manuscript must preserve that boundary.
\end{proof}

\begin{remark}
This screening statement is not a denial that those deeper or alternative manuscripts may matter strategically. It is a routing statement: they are not the default way to spend the next benchmark-facing move below the substrate-law frontier unless they do the same benchmark-facing work named above.
\end{remark}

\section{Conclusion}

The positive result of this note is routing discipline below the body-response-substrate-law frontier. The present theorem deriving the current observer-relevant pullback body-response core \(\PullbackBodyResponseCore\) from a benchmark-faithful section-centered displacement-response substrate law \(\SubstrateLaw\) is now the live benchmark-facing gain because it removes the current body-response core itself as the live unexplained stopping point on the recorded route.

The scope remains narrow and honest. The benchmark package is unchanged: the one operative metric is still exactly \(\CompletedMetric\), the ordinary matter route is unchanged, there is no second operative metric, and the low-energy matter law is not enlarged. Nothing here upgrades adoption to a completed final derivation or licenses stronger promotion language.

The next open burden is therefore clear. The next benchmark-facing manuscript should derive or justify the current displacement-response substrate law itself, or some nontrivial constituent or relation inside it, from still more primitive explicit substrate structure in a way that changes benchmark-facing status. Another same-law deeper relay, a re-expansion back to the larger already-spent presentations, or a side continuation that does not do that same benchmark-facing work is not the clean next manuscript target at current scope.

\end{document}
